% Appendix: Mathematical details for pairwise extremal dependence estimators
% This file is intended to be included in the thesis by \input{appendix_mathematics.tex}.
% Required packages in the main preamble:
% \usepackage{amsmath, amssymb, amsthm}

\appendix
\section{Mathematical details for pairwise extremal dependence estimators}\label{app:mathematics}

\subsection{Probability integral transform}\label{app:pit}

We prove that the marginal transform $U_{i,t}=F_i(X_{i,t})$ has a standard uniform distribution when $F_i$ is continuous.

\begin{proof}
    Let $X_{i,t}$ be a real-valued random variable with continuous distribution function $F_i$.
    Define
    \[
        U_{i,t}=F_i(X_{i,t}).
    \]
    For $u<0$,
    \[
        \Pr(U_{i,t}\leq u)=0.
    \]
    For $u>1$,
    \[
        \Pr(U_{i,t}\leq u)=1.
    \]
    Let $u\in[0,1]$.
    Define the generalized inverse
    \[
        F_i^{-1}(u)=\inf\{x\in\mathbb{R}:F_i(x)\geq u\}.
    \]
    Since $F_i$ is continuous,
    \[
        F_i(F_i^{-1}(u))=u.
    \]
    Then
    \begin{align*}
        \Pr(U_{i,t}\leq u)
         & =\Pr(F_i(X_{i,t})\leq u)      \\
         & =\Pr(X_{i,t}\leq F_i^{-1}(u)) \\
         & =F_i(F_i^{-1}(u))             \\
         & =u.
    \end{align*}
    Hence
    \[
        \Pr(U_{i,t}\leq u)=u,\qquad u\in[0,1].
    \]
    Therefore
    \[
        U_{i,t}\sim\mathrm{Uniform}(0,1).
    \]
    For two stations $i$ and $j$, define
    \[
        U_{i,t}=F_i(X_{i,t}),
        \qquad
        U_{j,t}=F_j(X_{j,t}).
    \]
    Then
    \[
        U_{i,t}\sim\mathrm{Uniform}(0,1),
        \qquad
        U_{j,t}\sim\mathrm{Uniform}(0,1).
    \]
    The marginal distributions are therefore equal after transformation, and the joint distribution of $(U_{i,t},U_{j,t})$ contains only dependence information.
\end{proof}

\subsection{Relation between $\chi_{ij}$ and $\theta_{ij}$ under max-stability}\label{app:chi-theta}

We prove that, under bivariate max-stability with standard uniform margins, the upper tail dependence coefficient and the extremal coefficient satisfy $\chi_{ij}=2-\theta_{ij}$.

\begin{proof}
    Let $(U_i,U_j)$ have standard uniform margins.
    Assume the bivariate extreme-value copula representation
    \[
        \Pr(U_i\leq u,U_j\leq u)=u^{\theta_{ij}},
        \qquad
        u\in(0,1),
    \]
    where
    \[
        \theta_{ij}\in[1,2].
    \]
    The upper tail dependence coefficient is
    \[
        \chi_{ij}
        =
        \lim_{u\uparrow 1}
        \Pr(U_j>u\mid U_i>u).
    \]
    By the definition of conditional probability,
    \begin{align*}
        \Pr(U_j>u\mid U_i>u)
         & =
        \frac{\Pr(U_i>u,U_j>u)}{\Pr(U_i>u)}.
    \end{align*}
    Since $U_i\sim\mathrm{Uniform}(0,1)$,
    \[
        \Pr(U_i>u)=1-u.
    \]
    By inclusion-exclusion,
    \begin{align*}
        \Pr(U_i>u,U_j>u)
         & =1-\Pr(U_i\leq u)-\Pr(U_j\leq u)+\Pr(U_i\leq u,U_j\leq u) \\
         & =1-u-u+u^{\theta_{ij}}                                    \\
         & =1-2u+u^{\theta_{ij}}.
    \end{align*}
    Therefore
    \begin{align*}
        \Pr(U_j>u\mid U_i>u)
         & =
        \frac{1-2u+u^{\theta_{ij}}}{1-u}.
    \end{align*}
    Thus
    \[
        \chi_{ij}
        =
        \lim_{u\uparrow 1}
        \frac{1-2u+u^{\theta_{ij}}}{1-u}.
    \]
    Both numerator and denominator converge to zero:
    \[
        \lim_{u\uparrow 1}(1-2u+u^{\theta_{ij}})=1-2+1=0,
    \]
    \[
        \lim_{u\uparrow 1}(1-u)=0.
    \]
    By l'Hopital's rule,
    \begin{align*}
        \chi_{ij}
         & =
        \lim_{u\uparrow 1}
        \frac{-2+\theta_{ij}u^{\theta_{ij}-1}}{-1} \\
         & =
        \frac{-2+\theta_{ij}}{-1}                  \\
         & =2-\theta_{ij}.
    \end{align*}
    Hence
    \[
        \boxed{\chi_{ij}=2-\theta_{ij}.}
    \]
    The endpoints are consistent with the interpretation:
    \[
        \theta_{ij}=1
        \quad\Longrightarrow\quad
        \chi_{ij}=1,
    \]
    \[
        \theta_{ij}=2
        \quad\Longrightarrow\quad
        \chi_{ij}=0.
    \]
\end{proof}

\subsection{Extremal coefficient as the exponent measure evaluated at $(1,1)$}\label{app:extremal-coefficient-exponent-measure}

We prove that, for a bivariate max-stable distribution with unit Fr\'echet margins, the pairwise extremal coefficient is $\theta_{ij}=V_{ij}(1,1)$.

\begin{proof}
    Let $(Z_i,Z_j)$ denote the bivariate vector
    \[
        (Z_i,Z_j)=(Z(s_i),Z(s_j)).
    \]
    Assume unit Fr\'echet margins:
    \[
        \Pr(Z_i\leq z)=\exp(-1/z),
        \qquad z>0,
    \]
    \[
        \Pr(Z_j\leq z)=\exp(-1/z),
        \qquad z>0.
    \]
    Let the joint distribution have exponent measure representation
    \[
        \Pr(Z_i\leq z_i,Z_j\leq z_j)
        =
        \exp\{-V_{ij}(z_i,z_j)\},
        \qquad z_i,z_j>0.
    \]
    Set
    \[
        z_i=z,
        \qquad
        z_j=z.
    \]
    Then
    \begin{align*}
        \Pr(Z_i\leq z,Z_j\leq z)
         & =\exp\{-V_{ij}(z,z)\}.
    \end{align*}
    For a max-stable distribution with unit Fr\'echet margins, the exponent measure is homogeneous of order $-1$:
    \[
        V_{ij}(az_i,az_j)=a^{-1}V_{ij}(z_i,z_j),
        \qquad a>0.
    \]
    Take
    \[
        a=z,
        \qquad
        z_i=1,
        \qquad
        z_j=1.
    \]
    Then
    \begin{align*}
        V_{ij}(z,z)
         & =V_{ij}(z\cdot 1,z\cdot 1) \\
         & =z^{-1}V_{ij}(1,1).
    \end{align*}
    Therefore
    \begin{align*}
        \Pr(Z_i\leq z,Z_j\leq z)
         & =\exp\{-z^{-1}V_{ij}(1,1)\}                 \\
         & =\exp\left\{-\frac{V_{ij}(1,1)}{z}\right\}.
    \end{align*}
    The pairwise extremal coefficient $\theta_{ij}$ is defined by
    \[
        \Pr(Z_i\leq z,Z_j\leq z)
        =
        \Pr(Z_i\leq z)^{\theta_{ij}}.
    \]
    Since
    \[
        \Pr(Z_i\leq z)=\exp(-1/z),
    \]
    we have
    \begin{align*}
        \Pr(Z_i\leq z)^{\theta_{ij}}
         & =\left[\exp(-1/z)\right]^{\theta_{ij}} \\
         & =\exp(-\theta_{ij}/z).
    \end{align*}
    Equating the two expressions for the same joint probability gives
    \[
        \exp\left\{-\frac{V_{ij}(1,1)}{z}\right\}
        =
        \exp\left\{-\frac{\theta_{ij}}{z}\right\}.
    \]
    Taking logarithms gives
    \[
        -\frac{V_{ij}(1,1)}{z}
        =
        -\frac{\theta_{ij}}{z}.
    \]
    Multiplying by $-z$ gives
    \[
        V_{ij}(1,1)=\theta_{ij}.
    \]
    Hence
    \[
        \boxed{\theta_{ij}=V_{ij}(1,1).}
    \]
    The bounds follow from the two limiting cases.
    Perfect dependence gives
    \[
        Z_i=Z_j,
    \]
    so
    \[
        \Pr(Z_i\leq z,Z_j\leq z)=\Pr(Z_i\leq z)=\exp(-1/z),
    \]
    which implies
    \[
        \theta_{ij}=1.
    \]
    Independence gives
    \begin{align*}
        \Pr(Z_i\leq z,Z_j\leq z)
         & =\Pr(Z_i\leq z)\Pr(Z_j\leq z) \\
         & =\exp(-1/z)\exp(-1/z)         \\
         & =\exp(-2/z),
    \end{align*}
    which implies
    \[
        \theta_{ij}=2.
    \]
    Therefore
    \[
        1\leq \theta_{ij}\leq 2.
    \]
\end{proof}

\subsection{F-madogram relation to the extremal coefficient}\label{app:f-madogram}

We prove that, under bivariate max-stability with standard uniform margins, the F-madogram satisfies $\theta_{ij}=(1+2\nu_F)/(1-2\nu_F)$.

\begin{proof}
    Let
    \[
        U_i=F_i(Z_i),
        \qquad
        U_j=F_j(Z_j),
    \]
    where $U_i$ and $U_j$ have standard uniform margins.
    The F-madogram is
    \[
        \nu_F(i,j)=\frac{1}{2}\mathbb{E}|U_i-U_j|.
    \]
    Use the identity
    \[
        |a-b|=a+b-2\min(a,b),
        \qquad a,b\in[0,1].
    \]
    Then
    \begin{align*}
        \mathbb{E}|U_i-U_j|
         & =\mathbb{E}(U_i+U_j-2\min(U_i,U_j))                    \\
         & =\mathbb{E}U_i+\mathbb{E}U_j-2\mathbb{E}\min(U_i,U_j).
    \end{align*}
    Since $U_i,U_j\sim\mathrm{Uniform}(0,1)$,
    \[
        \mathbb{E}U_i=\frac{1}{2},
        \qquad
        \mathbb{E}U_j=\frac{1}{2}.
    \]
    Therefore
    \begin{align*}
        \mathbb{E}|U_i-U_j|
         & =1-2\mathbb{E}\min(U_i,U_j).
    \end{align*}
    Thus
    \begin{align*}
        \nu_F(i,j)
         & =\frac{1}{2}-\mathbb{E}\min(U_i,U_j).
    \end{align*}
    It remains to compute $\mathbb{E}\min(U_i,U_j)$.
    For any random variable $Y\in[0,1]$,
    \[
        \mathbb{E}Y=\int_0^1 \Pr(Y>u)\,du.
    \]
    Set
    \[
        Y=\min(U_i,U_j).
    \]
    Then
    \begin{align*}
        \mathbb{E}\min(U_i,U_j)
         & =\int_0^1 \Pr(\min(U_i,U_j)>u)\,du \\
         & =\int_0^1 \Pr(U_i>u,U_j>u)\,du.
    \end{align*}
    By inclusion-exclusion,
    \begin{align*}
        \Pr(U_i>u,U_j>u)
         & =1-\Pr(U_i\leq u)-\Pr(U_j\leq u)+\Pr(U_i\leq u,U_j\leq u) \\
         & =1-u-u+\Pr(U_i\leq u,U_j\leq u)                           \\
         & =1-2u+\Pr(U_i\leq u,U_j\leq u).
    \end{align*}
    Under bivariate max-stability with extremal coefficient $\theta_{ij}$,
    \[
        \Pr(U_i\leq u,U_j\leq u)=u^{\theta_{ij}}.
    \]
    Therefore
    \begin{align*}
        \mathbb{E}\min(U_i,U_j)
         & =\int_0^1 \{1-2u+u^{\theta_{ij}}\}\,du                                                      \\
         & =\int_0^1 1\,du-2\int_0^1 u\,du+\int_0^1 u^{\theta_{ij}}\,du                                \\
         & =1-2\left[\frac{u^2}{2}\right]_0^1+\left[\frac{u^{\theta_{ij}+1}}{\theta_{ij}+1}\right]_0^1 \\
         & =1-1+\frac{1}{\theta_{ij}+1}                                                                \\
         & =\frac{1}{\theta_{ij}+1}.
    \end{align*}
    Hence
    \begin{align*}
        \nu_F(i,j)
         & =\frac{1}{2}-\frac{1}{\theta_{ij}+1}                               \\
         & =\frac{\theta_{ij}+1}{2(\theta_{ij}+1)}-\frac{2}{2(\theta_{ij}+1)} \\
         & =\frac{\theta_{ij}-1}{2(\theta_{ij}+1)}.
    \end{align*}
    Thus
    \[
        2\nu_F(i,j)=\frac{\theta_{ij}-1}{\theta_{ij}+1}.
    \]
    Solve for $\theta_{ij}$:
    \begin{align*}
        2\nu_F(\theta_{ij}+1)
         & =\theta_{ij}-1,             \\
        2\nu_F\theta_{ij}+2\nu_F
         & =\theta_{ij}-1,             \\
        \theta_{ij}-2\nu_F\theta_{ij}
         & =1+2\nu_F,                  \\
        \theta_{ij}(1-2\nu_F)
         & =1+2\nu_F,                  \\
        \theta_{ij}
         & =\frac{1+2\nu_F}{1-2\nu_F}.
    \end{align*}
    Hence
    \[
        \boxed{\theta_{ij}=\frac{1+2\nu_F(i,j)}{1-2\nu_F(i,j)}.}
    \]
    Replacing $\nu_F(i,j)$ by the empirical F-madogram
    \[
        \widehat{\nu}_F(i,j;c)
        =
        \frac{1}{2B^{(c)}}
        \sum_{b=1}^{B^{(c)}}
        \left|
        \widehat F_i^{(c)}(M_{i,b}^{(c)})
        -
        \widehat F_j^{(c)}(M_{j,b}^{(c)})
        \right|
    \]
    gives the estimator
    \[
        \widehat\theta_{ij}^{(c)}
        =
        \frac{1+2\widehat\nu_F(i,j;c)}{1-2\widehat\nu_F(i,j;c)}.
    \]
\end{proof}

\subsection{Consistency of the empirical finite-level tail-dependence estimator}\label{app:chi-consistency}

We prove that, under within-season stationarity and ergodicity, the empirical estimator $\widehat\chi_{ij}^{(c)}(u)$ converges to the finite-level conditional tail probability $\Pr(U_j^{(c)}>u\mid U_i^{(c)}>u)$.

\begin{proof}
    Fix a season
    \[
        c\in\{W,S\}.
    \]
    Fix a pair of stations
    \[
        (i,j).
    \]
    Let
    \[
        \{(U_{i,t}^{(c)},U_{j,t}^{(c)}):t\in T_c\}
    \]
    be the season-$c$ transformed pair process.
    Assume this process is strictly stationary and ergodic.
    Fix
    \[
        u\in(0,1).
    \]
    Define
    \[
        A_t^{(c)}(u)=\mathbf{1}\{U_{i,t}^{(c)}>u,\;U_{j,t}^{(c)}>u\},
    \]
    and
    \[
        B_t^{(c)}(u)=\mathbf{1}\{U_{i,t}^{(c)}>u\}.
    \]
    Since $(U_{i,t}^{(c)},U_{j,t}^{(c)})$ is strictly stationary and ergodic, any measurable transformation of it is also strictly stationary and ergodic.
    Therefore
    \[
        \{A_t^{(c)}(u)
        :
        t\in T_c\}
    \]
    is strictly stationary and ergodic, and
    \[
        \{B_t^{(c)}(u)
        :
        t\in T_c\}
    \]
    is strictly stationary and ergodic.
    Moreover,
    \[
        0\leq A_t^{(c)}(u)\leq 1,
        \qquad
        0\leq B_t^{(c)}(u)\leq 1.
    \]
    Hence
    \[
        \mathbb{E}|A_t^{(c)}(u)|<\infty,
        \qquad
        \mathbb{E}|B_t^{(c)}(u)|<\infty.
    \]
    By the ergodic theorem,
    \[
        \frac{1}{|T_c|}\sum_{t\in T_c} A_t^{(c)}(u)
        \xrightarrow{p}
        \mathbb{E}A_t^{(c)}(u),
    \]
    and
    \[
        \frac{1}{|T_c|}\sum_{t\in T_c} B_t^{(c)}(u)
        \xrightarrow{p}
        \mathbb{E}B_t^{(c)}(u).
    \]
    By the definition of $A_t^{(c)}(u)$,
    \begin{align*}
        \mathbb{E}A_t^{(c)}(u)
         & =\mathbb{E}\mathbf{1}\{U_{i,t}^{(c)}>u,\;U_{j,t}^{(c)}>u\} \\
         & =\Pr(U_{i,t}^{(c)}>u,\;U_{j,t}^{(c)}>u).
    \end{align*}
    By the definition of $B_t^{(c)}(u)$,
    \begin{align*}
        \mathbb{E}B_t^{(c)}(u)
         & =\mathbb{E}\mathbf{1}\{U_{i,t}^{(c)}>u\} \\
         & =\Pr(U_{i,t}^{(c)}>u).
    \end{align*}
    Since $U_{i,t}^{(c)}$ has a standard uniform margin,
    \[
        \Pr(U_{i,t}^{(c)}>u)=1-u.
    \]
    For $u\in(0,1)$,
    \[
        1-u>0.
    \]
    The empirical finite-level tail-dependence estimator is
    \[
        \widehat\chi_{ij}^{(c)}(u)
        =
        \frac{
            \sum_{t\in T_c}\mathbf{1}\{U_{i,t}^{(c)}>u,\;U_{j,t}^{(c)}>u\}
        }{
            \sum_{t\in T_c}\mathbf{1}\{U_{i,t}^{(c)}>u\}
        }.
    \]
    Using $A_t^{(c)}(u)$ and $B_t^{(c)}(u)$,
    \begin{align*}
        \widehat\chi_{ij}^{(c)}(u)
         & =
        \frac{\sum_{t\in T_c}A_t^{(c)}(u)}{\sum_{t\in T_c}B_t^{(c)}(u)} \\
         & =
        \frac{|T_c|^{-1}\sum_{t\in T_c}A_t^{(c)}(u)}{|T_c|^{-1}\sum_{t\in T_c}B_t^{(c)}(u)}.
    \end{align*}
    By the continuous mapping theorem,
    \begin{align*}
        \widehat\chi_{ij}^{(c)}(u)
         & \xrightarrow{p}
        \frac{\Pr(U_{i,t}^{(c)}>u,\;U_{j,t}^{(c)}>u)}{\Pr(U_{i,t}^{(c)}>u)} \\
         & =
        \Pr(U_{j,t}^{(c)}>u\mid U_{i,t}^{(c)}>u).
    \end{align*}
    Therefore
    \[
        \boxed{
            \widehat\chi_{ij}^{(c)}(u)
            \xrightarrow{p}
            \Pr(U_{j,t}^{(c)}>u\mid U_{i,t}^{(c)}>u)
        }
    \]
    for every fixed $u\in(0,1)$.
    If the limit
    \[
        \chi_{ij}^{(c)}
        =
        \lim_{u\uparrow 1}
        \Pr(U_{j,t}^{(c)}>u\mid U_{i,t}^{(c)}>u)
    \]
    exists, then the finite-level probability estimated above is the finite-$u$ approximation to the asymptotic tail-dependence coefficient.
\end{proof}
